\documentclass[conference]{IEEEtran}
\IEEEoverridecommandlockouts
\usepackage{cite}
\usepackage{amsmath,amssymb,amsthm}
\usepackage{graphicx}
\usepackage{booktabs}
\usepackage{multirow}
\usepackage[hidelinks]{hyperref}

% ---------------------------------------------------------------------------
% Audelle: measurement-grounded ML audio editing on commodity CPU hardware
% Every number in this paper is measured by the harness in Section VI.
% ---------------------------------------------------------------------------

\begin{document}
\title{Audelle: Measurement-Grounded Machine-Learning Audio\\ Editing on Commodity CPU Hardware}
\author{\IEEEauthorblockN{Lohit Aksha}}
\maketitle

\begin{abstract}
Audio editors increasingly ship ``AI'' features whose behaviour is asserted
rather than measured. We report Audelle, an editor in which every learned
component is accompanied by a quantitative claim, a ground-truth fixture that
supports it, and an explicit statement of what it cannot do. The system runs
entirely on a mobile CPU with no accelerator. It combines per-request model
selection for music source separation, statistical-model speech enhancement,
energy/spectral-flatness voice activity detection with clustering-based speaker
diarization, a timbre network trained on 3\,856 real instrument notes, a
phrase model learned from 1\,857 annotated melodic intervals, and
repetition-based audio inpainting. Natural-language requests are routed by a
deterministic grammar rather than a language model, so an edit cannot be
hallucinated. We measure the system against synthetic ground truth throughout:
distortion-to-noise ratio for speech enhancement, onset alignment error for
beat-locked layer placement, scale-invariant signal-to-distortion ratio for
inpainting, and per-clip accuracy for genre classification. Evaluation exposed
defects that code review had not, including a source separator that made speech
\emph{worse}, a voice activity detector that returned zero segments on
contiguous speech, and a melodic model whose regression head collapsed all
intervals to a single near-zero value. Two evaluation metrics were themselves
found to be invalid and were discarded rather than reported. A 99-case prompt
benchmark and 293 automated tests guard the resulting behaviour. We argue that
publishing negative results and invalidating one's own metrics is a necessary
part of deploying learned audio components in interactive tools.
\end{abstract}

\begin{IEEEkeywords}
music information retrieval, audio source separation, speech enhancement,
speaker diarization, audio inpainting, evaluation methodology, reproducible
research.
\end{IEEEkeywords}

\section{Introduction}

\subsection{The problem with unmeasured audio ML}

Modern audio editors advertise natural-language editing, stem extraction, voice
cleanup and ``seamless'' fill. In practice these features are frequently
demonstrations rather than mechanisms: a crossfade labelled as inpainting, a
pitch-swept sine labelled as a drum, a randomly populated bar chart labelled
as a stem analyser. Such interfaces are worse than no feature at all, because
they transfer the verification burden to the user.

Audelle is an audio editor whose design constraint is that no user-visible
behaviour may be shipped without a measurement behind it. Where a component
cannot be validated, it is relabelled or removed. This paper describes the
resulting system and, equally importantly, the evaluation methodology that made
it trustworthy, together with the defects that methodology uncovered.

\subsection{Contributions}

\begin{enumerate}
\item A CPU-only audio editing architecture that selects source-separation
      models per request rather than globally, on the principle that a
      music-trained model must not process speech (Section~III-B).
\item A timbre network trained from scratch on 3\,856 real notes, paired with a
      deterministic additive renderer, following the hybrid
      differentiablesignal-processing design (Section~III-E).
\item A phrase model that predicts \emph{where} rests fall, with the
      \emph{number} of rests supplied by a tempo-and-energy-aware arrangement
      model (Section~III-F).
\item Repetition-based inpainting that reconstructs a region and preserves the
      timeline, replacing a crossfade that deleted audio (Section~III-G).
\item A measurement methodology (Section~IV) built on synthesised ground-truth
      fixtures, and an account of the two evaluation metrics it invalidated and
      two of the six benchmark families it caught regressing.
\end{enumerate}

\subsection{Scope}

Audelle is an \emph{editing} tool, not a music generator. Large text-to-audio
generation was considered and rejected: it is the least differentiated component
against dedicated commercial products, the most expensive to run on the target
hardware, and orthogonal to Audelle's purpose. Where the system adds musical
material, that material is rendered by Audelle's own trained timbre model and
its own recorded drum bank.

\section{Related Work}

\subsection{Source separation}

Time--frequency masking with an ideal ratio mask remains a strong baseline, and
Conv-TasNet showed that learned time-domain masking can surpass it \cite{luo2019convtasnet}.
Spleeter packaged fixed pretrained models for fast separation \cite{hennequin2020spleeter}, and Open-Unmix
provided a reference implementation and dataset tooling for the same task \cite{stoter2019openunmix}.
Hybrid Transformer Demucs replaced the innermost bi-U-Net layers with a
cross-domain transformer using self-attention within a domain and cross-attention
across domains, reporting state-of-the-art results on MUSDB18 with additional
training data \cite{rouard2023hybrid}. Audelle uses Demucs, but its contribution here is a routing
policy: only \texttt{htdemucs\_6s} contains piano and guitar sources, so requests
naming those instruments cannot be served by the general model.

\subsection{Speech enhancement}

Statistical-model methods, in particular the Wiener formulation and its
decision-directed descendants, remain the reference against which learned
enhancers are measured \cite{ephraim1992statistical}. Audelle uses a decision-directed a-posteriori SNR
Wiener filter with noise-presence probability, selected because it requires no
training data and degrades gracefully on clean speech.

\subsection{Transcription and diarization}

Whisper demonstrated that large-scale weak supervision yields competitive
zero-shot recognition without fine-tuning \cite{radford2022whisper}. Voice activity detection by
adaptive energy thresholding with hangover, and clustering of MFCC statistics,
are long-established techniques \cite{davis1990comparison}; selection of the cluster count by the
silhouette coefficient \cite{rousseeuw1987silhouettes} follows Rousseeuw. Audelle implements this stack
directly rather than depending on a neural diarization pipeline, because the
target hardware has no accelerator and the pipeline must start in under a minute.

\subsection{Generative timbre}

Differentiable digital signal processing established that a synthesizer's
control parameters can be learned while its rendering stage remains hand-written
\cite{hayes2024ddsp}, and the design has since been reviewed and extended across music and speech
\cite{hayes2024ddsp}. Bazin et~al.\ cast note generation as spectrogram inpainting, using a
vector-quantised spectrogram autoencoder with a masked transformer decoder \cite{bazin2020spectrogram}.
Audelle adopts the hybrid philosophy of the former but not the generative model
of the latter: a small network predicts spectral envelopes and envelope
parameters, and an additive synthesizer renders audio, so no inference-time
vocoder is required.

\subsection{Repetition structure}

Rafii and Pardo's REPET exploited the repetitious structure of music to separate
vocals from accompaniment without a trained model \cite{rafii2013repet}. Audelle borrows the same
observation for a different task: if a region must be reconstructed, the best
available material is often the same musical position played elsewhere in the
track.

\subsection{Classification and features}

Mel-frequency cepstral coefficients remain a strong general-purpose acoustic
representation \cite{davis1990comparison}. Genre classification from spectral, chroma, energy and
tempo features is a standard benchmark \cite{tzanetakis2002genre}, and the random forest ensemble
remains a strong tabular baseline \cite{breiman2001random}. Audelle uses a permuted-temperament
key profile for key matching \cite{krumhansl1982perceived} and log-mel features for repetition search.
All signal processing is performed with \texttt{librosa} \cite{mcfee2015librosa}.

\section{System Architecture}

\subsection{Execution model}

Audelle is a Next.js front end with a FastAPI backend. Requests are resolved by
a two-stage planner. A deterministic grammar parses the utterance into an
\textit{intent} plus typed parameters; the executor then selects the algorithm.
Prompts that match no grammar rule are rejected with an explanation rather than
approximated. This is a deliberate design choice over a large language model: the
edit space is finite and enumerated, so a generative router can only introduce
failures. A latent language-model path is retained in the codebase and is inert
because no API key or local model is configured; it is documented as dormant
rather than deleted.

\subsection{Per-request separation model selection}

Separation is chosen by the \emph{requested source}, not globally. A general
four-stem model (\texttt{htdemucs\_ft}) serves vocals, drums and bass.
Piano and guitar requests route to \texttt{htdemucs\_6s}, the only available
model containing those stems. Source aliases are normalised through a resolver
that accepts natural phrasing; an instrument with no representable stem is
rejected explicitly rather than force-fitted, because silent coercion was the
cause of a shipped defect in which ``remove all instruments and keep the voice''
routed to vocal removal.

A second routing rule follows from domain match: speech is never processed by a
music-trained separator. Section~V-A reports the measurement that established
this rule, which reversed the system's original design.

\subsection{Statistical speech enhancement}

The denoiser estimates a smoothed noise floor over the shortest of the signal or
a 0.5\,s window, forms a decision-directed a-posteriori SNR estimate, and applies
a Wiener spectral gain with noise-presence probability, so that noise is
attenuated rather than eliminated during silent frames. A presence band lift
compensates for the perceived dullness left by broadband gain reduction. The
algorithm requires no training data and is invariant to the recording
microphone.

\subsection{Voice activity detection and diarization}

Frames are classified active when short-time energy exceeds an adaptive floor
defined in the log domain as
\begin{equation}
\tau = L_{\mathrm{noise}} + \max\!\left(5,\; 0.25\,(L_{90} - L_{10})\right),
\end{equation}
where $L_{10}$ and $L_{90}$ are the 10th and 90th percentiles of frame energy in
dB. Segments shorter than 0.5\,s are discarded and 0.25\,s of hangover is applied,
so that pauses inside a speaker turn do not split it. Each segment is summarised
by the mean and standard deviation of its MFCCs, and $k$-means partitions these;
$k$ is chosen by the silhouette coefficient. Adjacent segments assigned to the
same cluster are merged. Section~V-B reports why the original threshold returned
zero segments on contiguous speech.

\subsection{Learned timbre with a deterministic renderer}

A note is parameterised by a 32-band log-spectral envelope sampled over its
sustain, plus attack time, decay ratio, spectral centroid and spectral
flatness. The network \texttt{TimbreNet} maps an instrument embedding of
dimension 24, concatenated with normalised pitch and velocity, through two
hidden layers of 128 units to the 36-dimensional output. Partial amplitudes are
read off the predicted envelope with an inharmonicity term; flatness drives a
noise layer and the predicted parameters drive an ADSR envelope. Trained on
3\,856 notes spanning 27 instrument-family and source combinations, with
validation loss falling from 0.946 to 0.685, this is deliberately a hybrid: the
\emph{timbre is learned}, the renderer is ordinary additive synthesis, so a
wrong prediction produces a poor note rather than garbage. Instruments outside
the trained set return no prediction and fall back to a labelled DSP synth; no
fuzzy nearest-family match is attempted, for the same reason instrument
resolution is not.

\subsection{Phrase-level melodic model}

The phrase model learns the distribution of melodic intervals and durations
\emph{by classification into bins}, not by regression. Both the head
formulation and the arrangement stage are described below.

\paragraph{Classification, not regression.} An earlier head regressed the
interval and duration directly. A regressor trained under a conditional mean
loss returns the mean of its target distribution, which for a sparse interval
distribution is near zero; the measured mean absolute interval was 0.00
semitones and the generated lines were monotone. The model now classifies the
interval into 15 bins spanning $-7$ to $+7$ semitones, the duration into 16 bins
of which the zeroth denotes a rest, and predicts velocity directly, with the
architecture $5 \rightarrow 96 \rightarrow 96 \rightarrow 32$. Section~V-C
reports the before-and-after statistics.

\paragraph{Rests.} The model predicts the \emph{placement} of rests; the
arrangement stage decides the \emph{number}, by blending the model's rest
odds with a target fill rate as a geometric mean. This split avoids the
degenerate case in which the model, having learned that rests are rare,
emits none, and avoids the opposite failure in which rests are inserted at
random. Planning precedes rendering so that a note immediately repeated at the
same pitch is merged rather than re-struck, and an empty plan is guarded
against rather than rendered as silence.

\subsection{Rhythmic and tonal conditioning}

Detected tempo locks all generated layers: the drum sequencer places sixteenth
notes on the measured beat grid, and generated notes are snapped to the
detected beat phase. Detected key conditions melodic content. For two-file
jamming, keys are matched by correlating chroma histograms against
Krumhansl--Kessler key profiles \cite{krumhansl1982perceived}; alignment is refined by maximising
onset-envelope cross-correlation; and levels are matched by RMS.

\subsection{Recorded drum bank and rule-based arrangement}

Every drum hit Audelle plays is a recorded sample: 1\,200 one-shots across ten
percussion types, extracted and de-clicked at build time and stored as FLAC,
34.7\,MB in total. The previous implementation synthesised kicks from a
pitch-swept sine and snares from band-passed noise, which is why the ``add
drums'' feature sounded synthetic.

The distinction matters for what may be claimed. The \emph{samples} are real
recordings; the \emph{arrangement} --- which hits land on which sixteenth-note
step --- is a hand-written pattern library of fourteen styles, tempo-synced to
the measured beat grid. No trained drum model exists or is claimed. A test
verifies that the instrument catalogue advertised by the user interface matches
the instruments the backend can render, in both directions, so a style cannot be
promised and then found missing.

\subsection{Repetition-based inpainting}

Given a region to remove, the system seeks replacement material from the same
track, in three escalating strategies:

\begin{enumerate}
\item \textbf{Bar-aligned repeat.} If the region damages bar 9, the correct
      material is bar 5 or bar 13 --- the same musical position. The existing
      beat detector supplies the grid, so candidate positions one or more bars
      before the hole are enumerated directly, and each is scored by log-mel
      similarity to the audio on either side of the hole.
\item \textbf{Feature search.} A sliding search over the remaining track,
      scored identically, for material that is similar without being
      bar-aligned.
\item \textbf{Spectral fill.} Log-magnitude interpolation between the flanking
      frames, with phase propagated from the preceding frame's instantaneous
      frequency and corrected to land on the following frame.
\end{enumerate}

The timeline length is preserved, because a request to remove a transient
artifact means to keep the music, not to shorten the file. The metadata reports
which strategy ran, the source offset, and the match distance, so that a
fallback is never presented as a successful repeat.

\section{Measurement Methodology}

\subsection{Principle}

No claim is reported unless a measurement supports it, and every measurement
scores against ground truth whose value is known in advance. Synthetic
fixtures are constructed so that the correct answer is available by
construction: a loop built from a fixed period so that the phase-correct repeat
exists in the file; speech mixed at a known SNR so that the improvement is
computable; layer generated against a known grid so that alignment error is
exact.

Reconstruction quality is scored with the scale-invariant signal-to-distortion
ratio~\cite{leroux2019sdr}, chosen because it is invariant to the overall gain
mismatch between an estimate and its reference and therefore cannot be improved
merely by scaling the output --- a failure mode that plain SDR invites. Metrics
are chosen to be unfakeable: no metric in this paper rewards silence, since an
estimator that emits nothing would score infinitely well on several otherwise
tempting measures.

\subsection{Prompt benchmark}

A 99-case benchmark covers every intent the grammar supports, including
adversarial phrasings: ``phonk'' must not be captured by a rule matching
``honk''; ``cut pauses'' must not be captured by a trim rule; ``0.5 seconds''
must not be parsed as a duration from the wrong clause. All 99 cases pass.

\subsection{Automated regression suite}

293 automated tests execute. Beyond ordinary unit tests, the suite contains
ground-truth regression tests whose assertions are statements about measured
physics: that phrase statistics fall within measured bounds, that jam alignment
median error is below a threshold, that the filled region of an inpainted file
is not silent, and that the instrument catalogue advertised by the user
interface matches the instruments the backend can render.

\section{Results}

\subsection{Domain-matched separation (negative result)}

Routing speech through a music-trained separator \emph{degraded} quality. On
three broadcast-speech fixtures the SNR change ranged from $-4.3$\,dB to
$-7.2$\,dB. After replacing that path with a decision-directed Wiener filter,
the same fixtures measured $+1.3$, $+10.1$ and $+2.7$\,dB, with the voice
level raised 1--3\,dB and already-clean speech left essentially unchanged. This
result reversed the system's routing design. It is also the reason synthesised
speech was rejected as a test input: synthesised voices fall outside the
separator's training domain and yield near-silent stems, which would have
produced a spuriously excellent score.

\subsection{Adaptive VAD (negative result)}

The original threshold was $r > p_{40}(r)\cdot 1.5$, which fails whenever
speech occupies more than roughly 60\,\% of a file: the 40th percentile then
lies mid-speech and the multiplier excludes real speech, yielding
\emph{zero} segments. Equation~(1) removed this failure mode. Two further
defects surfaced only under test: cluster selection required at least $k+2$
segments, so a three-turn exchange could never resolve two speakers, and the
span-removal routine skipped any span beginning at sample~0, so leading audio
survived every ``keep only speaker 1'' request.

\subsection{Phrase model (quantitative)}

Across 13 annotated tracks yielding 1\,857 interval--duration pairs, the change
from regression to classification produced the statistics in
Table~\ref{tab:phrase}.

\begin{table}[t]
\centering
\caption{Measured phrase statistics, 13 tracks, 1\,857 pairs.}
\label{tab:phrase}
\begin{tabular}{lcc}
\toprule
Statistic & Regression & Classification \\
\midrule
Mean $|\Delta|$ (semitones)   & 0.00 & 4.69 \\
Distinct pitches per line    & 5    & 10 \\
Longest same-pitch run       & $15{+}$ & 2 \\
\bottomrule
\end{tabular}
\end{table}

A further defect was caught only by measuring the \emph{loudness} of quiet
tracks: a state slot conflated per-note amplitude with track energy, so quiet
inputs rendered as total silence. The same measurement revealed a fill density
of 10\,\%, against which the arrangement model was subsequently calibrated.

\subsection{Repetition-based inpainting}

On a fixture with a known 2\,s period and a hole punched into one repetition,
duration was preserved exactly ($8.00$\,s, against $7.50$\,s for the previous
splice-out) and the filled region was non-silent. We explicitly do not claim
that the original \emph{notes} are recovered: on through-composed material they
cannot be, because that audio no longer exists in the file. Seam continuity is
likewise not claimed as an improvement, since a 50\,ms equal-gain crossfade can
post a smaller single-sample step than the reconstructed splice; the tests
assert an upper bound rather than superiority. An intermediate design that
minimised seam discontinuity as its refinement objective was discarded after
measurement, because the objective is minimised at any quiet point and the
matcher consequently drifted away from the phase-correct repeat.

\subsection{Beat-locked drum placement}

Because all drum hits are now recordings, placement can be scored directly. The
energy of the onset-envelope restricted to the kick band was measured on-grid
against off-grid over generated patterns. On-grid placement yields $1.97\times$
the kick-band onset energy of off-grid placement, which confirms that the
sequencer is genuinely locked to the measured tempo rather than merely producing
something rhythmically plausible.

\subsection{Key-aware two-file fusion}

When a user supplies a second recording, the two are fused rather than merely
concatenated. Keys are matched by correlating chroma histograms against
Krumhansl--Kessler key profiles~\cite{krumhansl1982perceived}, so a track in A
minor is fused against one in B minor by transposition rather than by luck. The
beat phase is then aligned by maximising cross-correlation of the two
onset-envelope functions, and levels are matched by RMS. Alignment was measured
on synthetic fixtures of known tempo and key --- 120\,BPM A minor against
140\,BPM D minor --- and reaches a median onset-alignment error of 0.0\,ms with
a 90th percentile of 11.6\,ms, which is below the resolution at which the
discrepancy is audible as a flam.

\subsection{Genre classification}

The classifier is a random forest~\cite{breiman2001random} over 75 features
extracted per clip: the mean and standard deviation of 20 MFCCs and of a
12-bin chromagram, of spectral centroid, bandwidth, rolloff, zero-crossing rate
and RMS, together with estimated tempo. Training uses a balanced subset of 100
clips spanning the ten GTZAN genres~\cite{tzanetakis2002genre}. Results are
reported from a dedicated evaluation script with a fixed seed, which writes its
output to a JSON artefact in the repository; the earlier training script printed
its numbers to stdout and discarded the audio, so nothing on record stated what
the shipped artefact actually scored. Table~\ref{tab:genre} summarises the
result.

\begin{table}[t]
\centering
\caption{Genre classification, 100 GTZAN clips, 10 classes,
75 features, fixed seed. Chance is 10\,\%.}
\label{tab:genre}
\begin{tabular}{lcc}
\toprule
Protocol & Accuracy & Clips \\
\midrule
5-fold stratified CV & $75.0 \pm 8.9\,\%$ & 100 \\
30\,\% holdout        & $76.7\,\%$          & 30 \\
\midrule
Chance               & $10.0\,\%$          & --- \\
\bottomrule
\end{tabular}
\end{table}

Both protocols agree at roughly 7.5 times chance, which is consistent with the
literature on spectral and timbral genre cues and with the dataset size. The
per-fold spread of $60$--$85\,\%$ is wide relative to the mean, which is the
expected behaviour at this sample size and the reason the standard deviation is
reported rather than suppressed. Per-class holdout accuracy is in
\texttt{server/ml/models/genre\_eval.json}; it ranges from 100\,\% for
country, disco, metal and reggae to 33\,\% for pop, which is the class the
feature set distinguishes least. With 100 clips, no comparison against published
genre-classification results is meaningful, and none is made: the purpose of this
evaluation is to establish that the shipped artefact is better than chance and
to make that number reproducible, not to enter a benchmark.

\subsection{Prompt routing and regression}

All 99 benchmark cases pass and all 293 tests pass. The benchmark is the more
informative artefact: it encodes routing failures that unit tests would not
catch, including two cases where a shorter rule shadowed the correct one.

\subsection{Hardware envelope}

The entire pipeline runs on a mobile 8-core CPU. The timbre model trains in
roughly 30\,s and the drum bank builds in roughly 1\,min. Source separation is
the dominant cost of any request that invokes it.

\section{Threats to Validity}

The system was developed and evaluated on a single class of machine: a mobile
8-core CPU with 15.3\,GB of RAM and no GPU. All reported accuracies, energies
and alignments are therefore wall-clock-bound and are not claims about
throughput on accelerator hardware. Small-sample evaluation is a genuine
limitation: the genre classifier is trained on 100 clips, which is two orders of
magnitude smaller than a competitive benchmark, and the reported figures should
be read as a floor rather than an estimate of attainable accuracy. Ground-truth
fixtures are synthesised and therefore stylistically narrow; they are strong for
verifying numerical claims and weak for verifying perceptual ones. Finally, no
human listening test was conducted. Subjective quality claims are deliberately
absent from this paper.

\section{Failure Modes and Invalid Metrics}

We report two occasions on which our own evaluation was unsound, because both
would have produced favourable published numbers.

\paragraph{Synthetic speech against a music-trained separator.}
Evaluating separation on synthesised speech produced near-silent stems and an
excellent-looking score. Synthesised voices are outside the separator's training
domain; the metric was measuring the input, not the system. All separator
evaluation was moved to real recordings.

\paragraph{Comparing two different separation models against each other.}
Comparing the output of one separation model to the output of another produced
an apparent improvement of $-106$\,dB. The two models were being scored against
each other rather than against ground truth, and the number was an artefact of
mismatched assumptions. The comparison was discarded. After this, no metric is
reported without ground truth of known value.

\paragraph{Exceptions that hid errors.}
Two bare \texttt{except: pass} handlers silently swallowed genuine exceptions ---
a \texttt{NameError} in the timbre layer and a wrong keyword argument in the
denoiser --- after which the feature fell back to a default while still
reporting success. Failure is now surfaced in returned metadata, in the manner
of the existing \texttt{phrasing\_error} field.

\paragraph{Licensing.}
The drum bank is built from 1\,200 recorded one-shots obtained from a public
repository that does not declare a licence. We were unable to establish
permissive terms and regard the current redistribution as unresolved. This is
reported rather than silently shipped.

\section{Limitations and Future Work}

The phrase model is trained on 13 tracks, which is thin; the cache mechanism
avoids re-running separation, so scaling the corpus is inexpensive. The groove
classifier is currently trained on synthetic loops rather than real recordings,
which degrades the ``drums that match your song'' feature; retraining on real
audio is the next priority. Inpainting is verified for duration preservation and
for the absence of silence, but its perceptual quality is not established,
because ground-truth multitracks would be required. Inpainting cannot invent
music that was never played, and a learned generative model --- spectrogram
inpainting as in~\cite{bazin2020spectrogram}, or an audio language model --- would be required for that.
Audio-to-video muxing is implemented for extraction but not yet confirmed for
muxing. Precision sample-accurate waveform editing remains a gap relative to
workstation-class editors.

\section{Conclusion}

The principal claim of this work is methodological. Building an audio editor on
CPU-only hardware forces every learned component to be cheap, and cheapness
makes measurement easy to skip. When we measured anyway, the measurements found
a separator that made speech worse, a voice activity detector that returned
nothing on ordinary recordings, a melodic model that had collapsed to a
constant, an inpainter that deleted audio, and two evaluation metrics of our
own that were rewarding silence. A deterministic intent grammar, provenance in
returned metadata, and a benchmark that encodes routing as well as arithmetic
kept these from reaching users. Shipping a feature without a number behind it is
cheap and invisible; shipping the number is neither, and is the only way we
found to make the claim auditable.

\begin{thebibliography}{99}

\bibitem{luo2019convtasnet}
Y.~Luo and N.~Mesgarani, ``Conv-TasNet: surpassing ideal time-frequency
magnitude-based masking for speech separation,'' \emph{IEEE/ACM Trans. Audio,
Speech, Lang. Process.}, vol. 27, no. 8, pp. 1558--1571, 2019.
doi: 10.1109/TASLP.2019.2915167.

\bibitem{hennequin2020spleeter}
R.~Hennequin, A.~Khlif, F.~Voituret, and M.~Moussallam, ``Spleeter: a fast and
efficient music source separation tool with pre-trained models,''
\emph{J. Open Source Softw.}, vol. 5, no. 50, p. 2154, 2020.
doi: 10.21105/joss.02154.

\bibitem{stoter2019openunmix}
F.-R.~St{\'o}ter, S.~Uhlich, A.~Liutkus, and Y.~Mitsufuji, ``Open-unmix --- a
reference implementation for music source separation,'' \emph{J. Open Source
Softw.}, vol. 4, no. 37, p. 1667, 2019. doi: 10.21105/joss.01667.

\bibitem{rouard2023hybrid}
S.~Rouard, F.~Massa, and A.~D{\'e}fossez, ``Hybrid transformers for music
source separation,'' in \emph{Proc. IEEE Int. Conf. Acoust., Speech, Signal
Process. (ICASSP)}, 2023, pp. 1--5. doi: 10.1109/ICASSP49357.2023.10096956.

\bibitem{ephraim1992statistical}
Y.~Ephraim, ``Statistical-model-based speech enhancement systems,''
\emph{Proc. IEEE}, vol. 80, no. 10, pp. 1527--1558, 1992.
doi: 10.1109/5.168664.

\bibitem{radford2022whisper}
A.~Radford, J.~W.~Kim, T.~Xu, G.~Brockman, C.~McLeavey, and I.~Sutskever,
``Robust speech recognition via large-scale weak supervision,'' arXiv:2212.04356,
2022.

\bibitem{davis1990comparison}
S.~B.~Davis and P.~Mermelstein, ``Comparison of parametric representations for
monosyllabic word recognition in continuously spoken sentences,'' in
\emph{Readings in Speech Recognition}. Cambridge, MA, USA: MIT Press, 1990,
pp. 85--91. doi: 10.1016/B978-0-08-051584-7.50010-3.

\bibitem{rousseeuw1987silhouettes}
P.~J.~Rousseeuw, ``Silhouettes: a graphical aid to the interpretation and
validation of clustering analysis,'' \emph{J. Comput. Appl. Math.}, vol. 20,
pp. 53--65, 1987. doi: 10.1016/0377-0427(87)90125-7.

\bibitem{hayes2024ddsp}
B.~Hayes, J.~Shier, G.~Fazekas, A.~McPherson, and C.~Saitis, ``A review of
differentiable digital signal processing for music and speech synthesis,''
\emph{Front. Signal Process.}, vol. 2, art. no. 1284100, 2024.
doi: 10.3389/frsip.2023.1284100.

\bibitem{bazin2020spectrogram}
T.~Bazin, G.~Hadjeres, P.~Esling, and M.~Malt, ``Spectrogram inpainting for
interactive generation of instrument sounds,'' in \emph{Proc. Int. Conf. on AI
Music Creativity (AIMC)}, Stockholm, Sweden, 2020, pp. 306--312.
doi: 10.5281/zenodo.4285406.

\bibitem{rafii2013repet}
Z.~Rafii and B.~Pardo, ``REpeating pattern extraction technique (REPET): a
simple method for music/voice separation,'' \emph{IEEE/ACM Trans. Audio,
Speech, Lang. Process.}, vol. 21, no. 1, pp. 73--84, 2013.
doi: 10.1109/TASL.2012.2213249.

\bibitem{tzanetakis2002genre}
G.~Tzanetakis and P.~Cook, ``Musical genre classification of audio signals,''
\emph{IEEE Trans. Speech Audio Process.}, vol. 10, no. 5, pp. 302--303, 2002.
doi: 10.1109/TSA.2002.800560.

\bibitem{breiman2001random}
L.~Breiman, ``Random forests,'' \emph{Mach. Learn.}, vol. 45, no. 1, pp. 5--32,
2001. doi: 10.1023/A:1010933404324.

\bibitem{krumhansl1982perceived}
C.~L.~Krumhansl, J.~J.~Bharucha, and E.~J.~Kessler, ``Perceived harmonic
structure of chords in three related musical keys,'' \emph{J. Exp. Psychol.:
Hum. Percept. Perform.}, vol. 8, no. 1, pp. 24--31, 1982.
doi: 10.1037/0096-1523.8.1.24.

\bibitem{mcfee2015librosa}
B.~McFee, C.~Raffel, D.~Liang, D.~Ellis, M.~McVicar, E.~Battenberg, and
O.~Nieto, ``librosa: audio and music signal analysis in Python,'' in
\emph{Proc. 14th Python in Science Conf. (SciPy)}, 2015, pp. 18--25.
doi: 10.25080/Majora-7b98e3ed-003.

\bibitem{leroux2019sdr}
J.~Le Roux, S.~Wisdom, H.~Erdogan, and J.~R.~Hershey, ``SDR --- half-baked or
well done?'' in \emph{Proc. IEEE Int. Conf. Acoust., Speech, Signal Process.
(ICASSP)}, 2019, pp. 626--630. doi: 10.1109/ICASSP.2019.8683855.

\end{thebibliography}

\end{document}